%%%%%%%%%%%%%%%%%%%%%%%%%%%%%%%%%%%%%%%%%%
% Short Sectioned Assignment LaTeX Template Version 1.0 (5/5/12)
% This template has been downloaded from: http://www.LaTeXTemplates.com
% Original author:  Frits Wenneker (http://www.howtotex.com)
% License: CC BY-NC-SA 3.0 (http://creativecommons.org/licenses/by-nc-sa/3.0/)
%%%%%%%%%%%%%%%%%%%%%%%%%%%%%%%%%%%%%%%%%

%----------------------------------------------------------------------------------------
%	PACKAGES AND OTHER DOCUMENT CONFIGURATIONS
%----------------------------------------------------------------------------------------

\documentclass[paper=a4, fontsize=11pt]{scrartcl} % A4 paper and 11pt font size

% ---- Entrada y salida de texto -----

\usepackage[T1]{fontenc} % Use 8-bit encoding that has 256 glyphs
\usepackage[utf8]{inputenc}
%\usepackage{fourier} % Use the Adobe Utopia font for the document - comment this line to return to the LaTeX default

% ---- Idioma --------

\usepackage[spanish, es-tabla]{babel} % Selecciona el español para palabras introducidas automáticamente, p.ej. "septiembre" en la fecha y especifica que se use la palabra Tabla en vez de Cuadro

% ---- Otros paquetes ----

\usepackage{url} % ,href} %para incluir URLs e hipervínculos dentro del texto (aunque hay que instalar href)
\usepackage{amsmath,amsfonts,amsthm} % Math packages
%\usepackage{graphics,graphicx, floatrow} %para incluir imágenes y notas en las imágenes
\usepackage{graphics,graphicx, float} %para incluir imágenes y colocarlas

% Para hacer tablas comlejas
%\usepackage{multirow}
%\usepackage{threeparttable}

%\usepackage{sectsty} % Allows customizing section commands
%\allsectionsfont{\centering \normalfont\scshape} % Make all sections centered, the default font and small caps

\usepackage{fancyhdr} % Custom headers and footers
\pagestyle{fancyplain} % Makes all pages in the document conform to the custom headers and footers
\fancyhead{} % No page header - if you want one, create it in the same way as the footers below
\fancyfoot[L]{} % Empty left footer
\fancyfoot[C]{} % Empty center footer
\fancyfoot[R]{\thepage} % Page numbering for right footer
\renewcommand{\headrulewidth}{0pt} % Remove header underlines
\renewcommand{\footrulewidth}{0pt} % Remove footer underlines
\setlength{\headheight}{13.6pt} % Customize the height of the header

\numberwithin{equation}{section} % Number equations within sections (i.e. 1.1, 1.2, 2.1, 2.2 instead of 1, 2, 3, 4)
\numberwithin{figure}{section} % Number figures within sections (i.e. 1.1, 1.2, 2.1, 2.2 instead of 1, 2, 3, 4)
\numberwithin{table}{section} % Number tables within sections (i.e. 1.1, 1.2, 2.1, 2.2 instead of 1, 2, 3, 4)

\setlength\parindent{0pt} % Removes all indentation from paragraphs - comment this line for an assignment with lots of text

\newcommand{\horrule}[1]{\rule{\linewidth}{#1}} % Create horizontal rule command with 1 argument of height






%----------------------------------------------------------------------------------------
%	TÍTULO Y DATOS DEL ALUMNO
%----------------------------------------------------------------------------------------

\title{	
\normalfont \normalsize 
\textsc{\textbf{Grundlagen der Wissensverarbeitung} \\ Computer Science \\ Universität Hamburg} \\ [25pt] % Your university, school and/or department name(s)
~\\
~\\
~\\
\horrule{0.5pt} \\[0.4cm] % Thin top horizontal rule
\Huge Tutorial 7: Constraint Satisfaction\\ % The assignment title
\horrule{2pt} \\[0.5cm] % Thick bottom horizontal rule
~\\
~\\
}

\author{Rafael Ruz Gómez\\Miguel Robles Urquiza} % Nombre y apellidos

\date{\normalsize \today} % Incluye la fecha actual

%----------------------------------------------------------------------------------------
% DOCUMENTO
%----------------------------------------------------------------------------------------

\begin{document}

\maketitle % Muestra el Título

\begin{figure}
	\centering
	\includegraphics[scale=0.8]{logo_uni_hamburg.png}
\end{figure}

\newpage %inserta un salto de página




%----------------------------------------------------------------------------------------
%	Question 1
%----------------------------------------------------------------------------------------

\section*{Exercise 7.1}

\subsection*{Formalize this riddle in the form of a constraint network, with the constraints being reasonably small. (I.e. writing a single constraint is not a good solution!) In the following pattern each letter stands for a digit so that the resulting sum is correct.}
\textbf{
\ \ \ \ S E N D\\
+\  M O R E\\
=======\\
\ M O N E Y\\
}

Let's call $S_i$ the sum done in the last step. Then the constraints are:
\begin{itemize}
	\item Different letter have different values: $S \neq E , E \neq N , N \neq D...$
	\item $(D + E) \% 10 == Y \ \ \ \ \ \ ( S_0 = (D + E) )$
	\item $ (((S_0 - S_0) \% 10) / 10 + N + R) \% 10 == E \ \ \ \ \ \ ( S_1 = ((S_0 - S_0) \% 10) / 10 + N + R ) $
	\item $ (((S_1 - S_1) \% 10) / 10 + E + O) \% 10 == N \ \ \ \ \ \ ( S_2 = ((S_1 - S_1) \% 10) / 10 + E + O) $
	\item $ (((S_2 - S_2) \% 10) / 10 + S + M) \% 10 == O \ \ \ \ \ \ ( S_3 = ((S_2 - S_2) \% 10) / 10 + S + M) $
	\item $ (((S_3 - S_3) \% 10) / 10 == M $
\end{itemize}
%----------------------------------------------------------------------------------------
%	Question 2
%----------------------------------------------------------------------------------------

\subsection*{Manual constraint solving.}

\textbf{
Crossword puzzles are often used in newspapers because they provide joy in solving semi-complex problems by combining logics and human experience. For the crossword above we want to find 6 words of length 3 that fit into the 3x3 table in a way that 3 words can be read horizontal from left to right and 3 words can be read vertically from top to bottom. Choose the words from the following list:\\
\\
add, ado, age, ago, aid, ail, aim, air,\\
and, any, ape, apt, arc, are, ark, arm,\\
art, ash, ask, auk, awe, awl, aye, bad,\\
bag, ban, bat, bee, boa, ear, eel, eft,\\
far, fat, fit, lee, oaf, rat, tar, tie.\\
\\
First, try to solve the problem without any formal methods or tools. How do you approach this problem as a human? (It is not necessary to give a full solution to the problem at this point, but you should report on the strategies you employ as a human and the problems you encounter.)
}
\newpage

A good way to begin is seeing which letter appearing more. We take one of this words, let's think the word is "XYZ". \\

After that, it's time to check if there is other word that starts with "X". If that happens, we should put the word in the first row or in the first colum and add the other word in the correct position.\\
If don't, we shoul check if there is other word that starts with "Y". If that happens, we should put the word in the first row or in the first colum and add the other word in the correct position.\\
 If don't, we should check if there is other word that starts with "Z". If that happens, we should put the word in the first row or in the first colum and add the other word in the correct position. \\
  If don't, we should check if there is other word that has a "X" \ in second position. If that happens, we should put the word in the second row or second coulm and add the other word in the correct position. \\
 If don't, we should check if there is other word that has a "Y" \ in second position. If that happens, we should put the word in the second row or second coulm and add the other word in the correct position. \\
 If don't, we should check if there is other word that has a "Z" \ in second position. If thath happens, we should put the word in the second row or second coulm.\\
  If don't, we should check if there is other word that has a "X" \ in third position. If thath happens, we should put the word in the third row or third coulm and add the other word in the correct position.\\
 If don't, we should check if there is other word that has a "Y" \ in third position. If thath happens, we should put the word in the third row or third coulm and add the other word in the correct position.\\
  If don't, we should check if there is other word that has a "Z" \ in third position. If thath happens, we should put the word in the third row or third coulm and add the other word in the correct position.\\
  
Basically, the constraint is that the letters where a pair of words intersect must be the same.
 


%----------------------------------------------------------------------------------------
%	Question 3
%----------------------------------------------------------------------------------------

\subsection*{Solve the problem by hand using domain consistency as a first step and as a second step the arc consistency. Document this process thoroughly.}

 \begin{itemize}
 
 \item Domain consistency: 
 
We can solve the problem by treating it as a CSP. The set of variables is \newline S = "X,Y,Z", that represents X as the first letter of the word, Y as the second and Z as the last. Dom(X) = Dom(Y) = Dom(Z) (I.e All the letters of the alphabet)\\

Now, we take the first letter of the alphabet and check if it's appears in any of the words given. If not, we delete that letter from the domain.



 
 
 

 \end{itemize}

%----------------------------------------------------------------------------------------
%	Question 4
%----------------------------------------------------------------------------------------


\subsection*{• Implement the arc consistency algorithm (found in sect. 4.5 of Poole and Mackworth (2010)) along with a suitable representation of the problem to solve this puzzle.}






%------------------------------------------------

%\bibliography{citas} %archivo citas.bib que contiene las entradas 
%\bibliographystyle{plain} % hay varias formas de citar

\end{document}

